% ============================================================
% UNIDAD: DE LA IDEA AL PLAN DE NEGOCIOS
% ============================================================
\clearpage
\section*{UNIDAD: DE LA IDEA AL PLAN DE NEGOCIOS}
\subsection*{Laboratorio de ideas, factibilidad y comunicación}

\begin{recuadro}
\textbf{Objetivo de la unidad:} transformar una idea en un proyecto evaluado: detectar
oportunidades, ordenar el modelo de negocio, calcular si es económicamente viable y
presentarlo de forma clara y convincente.
\end{recuadro}

% ------------------------------------------------------------
\section*{1. LABORATORIO DE IDEAS Y OPORTUNIDADES}

Una buena idea de negocio empieza por un \textbf{problema real} de alguien que esté dispuesto
a pagar por resolverlo.

\subsubsection*{1.1. Fuentes de oportunidades}
\begin{itemize}
  \item Problemas que vemos en el taller, en el trabajo o en el barrio.
  \item Cambios tecnológicos (drones, ADS-B, electrificación, IA) y normativos (nuevas
        exigencias que obligan a modernizar equipos).
  \item Servicios que hoy se importan o que no existen en la zona.
  \item Mejorar algo que ya existe: más barato, más rápido, más cerca, mejor atención.
\end{itemize}

\subsubsection*{1.2. Técnicas de generación de ideas}
\begin{itemize}
  \item \textbf{Lluvia de ideas:} cantidad antes que calidad, sin criticar en la primera ronda.
  \item \textbf{Design thinking:} empatizar con el usuario, definir el problema, idear,
        construir un prototipo y probarlo.
  \item \textbf{Validación:} antes de invertir, hablar con 10 o 20 posibles clientes y
        preguntar por el problema (no por la idea). Construir un \textbf{producto mínimo
        viable} (MVP) para probar con poco dinero.
\end{itemize}

% ------------------------------------------------------------
\section*{2. EL MODELO DE NEGOCIO: LIENZO CANVAS}

El \textbf{Business Model Canvas} resume en una sola hoja cómo un emprendimiento crea,
entrega y captura valor. Tiene nueve bloques:

\begin{center}
\small
\begin{tabular}{|p{0.25\linewidth}|p{0.65\linewidth}|}
\hline
\textbf{Bloque} & \textbf{Pregunta} \\ \hline
1. Segmentos de clientes & ¿Para quién creamos valor? \\ \hline
2. Propuesta de valor & ¿Qué problema resolvemos? ¿Por qué nos elegirían? \\ \hline
3. Canales & ¿Cómo llegamos al cliente y le entregamos el servicio? \\ \hline
4. Relación con clientes & ¿Qué tipo de vínculo mantenemos (personal, digital, contrato)? \\ \hline
5. Fuentes de ingresos & ¿Por qué y cómo nos pagan? \\ \hline
6. Recursos clave & ¿Qué necesitamos tener (habilitaciones, instrumental, personal)? \\ \hline
7. Actividades clave & ¿Qué tenemos que hacer bien? \\ \hline
8. Socios clave & ¿Quiénes nos ayudan (proveedores, representantes, talleres)? \\ \hline
9. Estructura de costos & ¿Cuáles son los costos más importantes? \\ \hline
\end{tabular}
\end{center}

% ------------------------------------------------------------
\section*{3. EL PLAN DE NEGOCIOS}

El \textbf{plan de negocios} es el documento que desarrolla el proyecto en detalle. Sirve para
pensar, para decidir y para conseguir socios, créditos o inversores.

\begin{enumerate}
  \item \textbf{Resumen ejecutivo:} el proyecto en una página (se escribe al final).
  \item \textbf{Descripción del proyecto y del equipo.}
  \item \textbf{Análisis del mercado:} clientes, tamaño del mercado, competidores y precios.
  \item \textbf{Análisis FODA} (ver unidad de Calidad).
  \item \textbf{Plan de marketing:} producto, precio, plaza (canales) y promoción.
  \item \textbf{Plan de operaciones:} procesos, instalaciones, equipamiento, proveedores,
        habilitaciones y normativa (en aviación, certificación ANAC del taller).
  \item \textbf{Organización y recursos humanos.}
  \item \textbf{Forma jurídica e impuestos.}
  \item \textbf{Plan económico-financiero:} inversión, costos, ventas, punto de equilibrio,
        flujo de fondos y evaluación (VAN, TIR, recupero).
  \item \textbf{Riesgos y plan de contingencia.}
  \item \textbf{Impacto social y ambiental.}
\end{enumerate}

% ------------------------------------------------------------
\section*{4. FACTIBILIDAD TÉCNICO-ECONÓMICA}

\subsubsection*{4.1. Costos fijos y variables}
\begin{itemize}
  \item \textbf{Costos fijos:} no cambian con la cantidad producida (alquiler, sueldos,
        seguros, contador, internet).
  \item \textbf{Costos variables:} aumentan con cada unidad vendida (componentes, insumos,
        envíos, comisiones).
  \item \textbf{Contribución marginal (CM):} precio de venta $-$ costo variable unitario. Es
        lo que cada unidad aporta para cubrir los costos fijos.
\end{itemize}

\subsubsection*{4.2. Punto de equilibrio}
Es la cantidad que hay que vender para \textbf{no ganar ni perder}:
\[
Q_{\mbox{equilibrio}} = \frac{\mbox{Costos fijos}}{\mbox{Precio} - \mbox{Costo variable unitario}}
\]

\subsubsection*{4.3. Flujo de fondos y evaluación}
\begin{itemize}
  \item \textbf{Flujo de fondos:} ingresos menos egresos de dinero de cada período, empezando
        por la inversión inicial (con signo negativo).
  \item \textbf{Período de recupero:} tiempo que tarda el proyecto en devolver la inversión.
  \item \textbf{VAN (Valor Actual Neto):} suma de los flujos futuros ``traídos al presente''
        con una tasa de descuento (lo que podríamos ganar en otra alternativa), menos la
        inversión. Si el \textbf{VAN es positivo}, el proyecto conviene.
        \[
        VAN = -I_0 + \sum_{t=1}^{n} \frac{F_t}{(1+i)^t}
        \]
  \item \textbf{TIR (Tasa Interna de Retorno):} la tasa que hace VAN = 0. Si la TIR es mayor
        que la tasa de descuento, el proyecto conviene.
\end{itemize}

\begin{recuadro}
\textbf{Ejemplo: taller de reparación de equipos electrónicos de aviónica.}

\textit{Datos:} inversión inicial \$12.000.000 (instrumental, herramientas y adecuación del
local). Precio promedio por reparación \$150.000; costo variable (componentes e insumos)
\$50.000. Costos fijos mensuales \$2.500.000 (alquiler, servicios, contador e ingreso del
técnico).

\textit{Contribución marginal:} \$150.000 $-$ \$50.000 = \$100.000 por reparación.

\textit{Punto de equilibrio:} \$2.500.000 / \$100.000 = \textbf{25 reparaciones por mes}.

\textit{Flujo de fondos} (tasa de descuento 20\% anual en términos reales): año 1:
\$8.000.000 (arranque); años 2 a 5: \$12.000.000 por año (35 reparaciones/mes).

\textit{Recupero:} al terminar el año 1 se recuperaron \$8.000.000; faltan \$4.000.000, que
se recuperan en el primer tercio del año 2 $\Rightarrow$ \textbf{aprox. 1 año y 4 meses}.

\textit{VAN:} $-12.000.000 + \frac{8.000.000}{1,2} + \frac{12.000.000}{1,2^2} +
\frac{12.000.000}{1,2^3} + \frac{12.000.000}{1,2^4} + \frac{12.000.000}{1,2^5}
\approx$ \textbf{\$20.554.000} $> 0$ $\Rightarrow$ conviene.

\textit{TIR:} \textbf{aprox. 80\% anual}, muy superior a la tasa de descuento.
\end{recuadro}

\paragraph*{Factibilidad técnica y ambiental.}
Además de los números, hay que verificar que el proyecto se pueda hacer: habilitaciones (en
aviación, un taller debe estar certificado por la ANAC bajo RAAC 145), personal con licencia,
instrumental calibrado, proveedores confiables, y que cumpla la normativa ambiental (gestión
de residuos, efluentes, ruido).

% ------------------------------------------------------------
\section*{5. TÉCNICAS DE COMUNICACIÓN}

\subsubsection*{5.1. El pitch}
Un \textbf{pitch} es una presentación breve (de 1 a 5 minutos) para convencer a un inversor,
un jurado, un cliente o un socio. Estructura recomendada:
\begin{enumerate}
  \item \textbf{Gancho:} una frase, un dato o una pregunta que capte la atención.
  \item \textbf{Problema:} qué duele y a quién.
  \item \textbf{Solución:} nuestra propuesta y por qué es mejor.
  \item \textbf{Mercado:} cuántos clientes hay y cuánto pagan.
  \item \textbf{Modelo de negocio:} cómo ganamos dinero.
  \item \textbf{Equipo:} por qué nosotros podemos hacerlo.
  \item \textbf{Pedido:} qué necesitamos (inversión, socio, prueba piloto).
\end{enumerate}

\subsubsection*{5.2. Hablar en público}
\begin{itemize}
  \item Ensayar en voz alta y medir el tiempo.
  \item Pocas diapositivas, con imágenes y pocas palabras; nunca leer la pantalla.
  \item Mirar al público, hablar despacio, usar ejemplos concretos.
  \item Preparar las preguntas difíciles (costos, competencia, riesgos).
\end{itemize}

\subsubsection*{5.3. Comunicación escrita profesional}
Un correo formal tiene: asunto claro, saludo, motivo en la primera línea, desarrollo breve,
pedido concreto, despedida y firma con nombre, cargo y teléfono. Revisar la ortografía y los
adjuntos antes de enviar.

% ------------------------------------------------------------
\section*{6. ACTIVIDADES DE LA UNIDAD}

\begin{enumerate}
  \item \textbf{Laboratorio de ideas:} en grupos, listar 10 problemas del sector aeronáutico o
        de la comunidad y elegir uno para desarrollar.
  \item \textbf{Canvas:} completar el lienzo del proyecto final del grupo.
  \item \textbf{Factibilidad:} calcular el punto de equilibrio, el flujo de fondos a 5 años,
        el VAN, la TIR y el recupero del proyecto, en una planilla de cálculo.
  \item \textbf{Plan de negocios:} redactar el plan completo siguiendo la estructura del
        punto 3.
  \item \textbf{Pitch final:} presentar el proyecto en 3 minutos ante un jurado (docentes,
        invitados de empresas o de la agencia de desarrollo local).
\end{enumerate}
